\item Une entreprise cherche à expliquer le salaire annuel à l'aide du nombre
d'années d'expérience $x_1$, du nombre de mois d'expérience $x_2$ et du niveau
d'études $x_3$.

\begin{enumerate}[label=\alph*)]
  \item Expliquez pourquoi $\bX\trans\bX$ ne sera pas inversible si
        $x_2=12x_1$ pour chaque employé.
  \item Quelle modification simple du modèle proposez-vous?
  \item Donnez un autre exemple de variables explicatives qui produiraient une
        dépendance linéaire exacte entre les colonnes de $\bX$.
  \item Le fait d'avoir deux variables fortement corrélées, mais non
        parfaitement corrélées, rend-il nécessairement
        $\bX\trans\bX$ non inversible?
\end{enumerate}

\prof{
\begin{enumerate}[label=\alph*)]
  \item La colonne de $\bX$ associée à $x_2$ est exactement 12 fois la colonne
        associée à $x_1$. Les colonnes de $\bX$ sont donc linéairement
        dépendantes. Ainsi, $\bX$ n'est pas de plein rang colonne et
        $\bX\trans\bX$ est singulière.

  \item Il suffit de retirer l'une des deux variables $x_1$ et $x_2$, puisqu'elles
        contiennent exactement la même information exprimée dans des unités
        différentes. Le choix de l'unité peut être guidé par la facilité
        d'interprétation.

  \item Par exemple, si un modèle contient simultanément la température en
        degrés Celsius $x_C$, la température en degrés Fahrenheit $x_F$ et une
        ordonnée à l'origine, alors
        \[
          x_F=32\mathbf{1}+\frac{9}{5}x_C.
        \]
        La colonne de $x_F$ est une combinaison linéaire de la colonne de 1 et
        de celle de $x_C$.

  \item Non. Une forte corrélation ne constitue pas une dépendance linéaire
        exacte. La matrice peut demeurer inversible, mais être proche d'une
        matrice singulière. Les estimations peuvent alors être instables et
        leurs variances élevées : c'est un problème de multicolinéarité.
\end{enumerate}
}